% =============================================================================
%  pretextual/listas_siglas.tex
%  Contém: Lista de Figuras, Lista de Tabelas, Lista de Códigos e
%          Lista de Siglas e Abreviaturas
%
%  As três primeiras são AUTOMÁTICAS — geradas a partir dos \caption{} do texto.
%  A lista de siglas é MANUAL — você adiciona cada entrada aqui.
%
%  QUANDO INCLUIR CADA LISTA (recomendação ABNT):
%    - Lista de Figuras  → se o trabalho tiver 5 ou mais figuras
%    - Lista de Tabelas  → se o trabalho tiver 5 ou mais tabelas
%    - Lista de Códigos  → se houver vários blocos de código com caption
%    - Lista de Siglas   → se o texto usar muitas siglas técnicas
%  Se não quiser uma lista, basta comentar o bloco correspondente.
% =============================================================================


% -----------------------------------------------------------------------------
%  LISTA DE FIGURAS
%  Gerada automaticamente. O asterisco (*) impede que a lista apareça
%  no próprio sumário (comportamento correto conforme ABNT).
% -----------------------------------------------------------------------------
\pdfbookmark[0]{\listfigurename}{lof}  % Adiciona bookmark no PDF
\listoffigures*
\cleardoublepage  % Garante que o próximo elemento comece em página nova


% -----------------------------------------------------------------------------
%  LISTA DE TABELAS
% -----------------------------------------------------------------------------
\pdfbookmark[0]{\listtablename}{lot}
\listoftables*
\cleardoublepage


% -----------------------------------------------------------------------------
%  LISTA DE CÓDIGOS
%  Exibe todos os blocos \begin{lstlisting}[caption={...}] do trabalho.
%  Remova este bloco se não quiser listar os trechos de código separadamente.
% -----------------------------------------------------------------------------
\pdfbookmark[0]{Lista de Códigos}{lol}
\lstlistoflistings
\cleardoublepage


% -----------------------------------------------------------------------------
%  LISTA DE SIGLAS E ABREVIATURAS
%  Feita MANUALMENTE. Mantenha em ordem ALFABÉTICA.
%
%  Formato: \item[SIGLA]  Significado por extenso
%
%  Alternativa automática: pacote 'glossaries' (descomente no main.tex)
%  e defina cada sigla com: \newacronym{chave}{SIGLA}{Significado}
% -----------------------------------------------------------------------------
\begin{siglas}
  \item[ABNT]   Associação Brasileira de Normas Técnicas
  \item[API]    Application Programming Interface
  \item[BD]     Banco de Dados
  \item[CRUD]   Create, Read, Update, Delete
  \item[CSS]    Cascading Style Sheets
  \item[DER]    Diagrama Entidade-Relacionamento
  \item[HTML]   HyperText Markup Language
  \item[HTTP]   HyperText Transfer Protocol
  \item[HTTPS]  HyperText Transfer Protocol Secure
  \item[IDE]    Integrated Development Environment
  \item[JSON]   JavaScript Object Notation
  \item[MVC]    Model-View-Controller
  \item[RF]     Requisito Funcional
  \item[RNF]    Requisito Não Funcional
  \item[REST]   Representational State Transfer
  \item[SQL]    Structured Query Language
  \item[TCC]    Trabalho de Conclusão de Curso
  \item[UML]    Unified Modeling Language
  \item[URL]    Uniform Resource Locator
  % TODO: Adicione as siglas específicas do seu projeto abaixo:
\end{siglas}
